\documentclass[11pt]{article}

\usepackage[T1]{fontenc}
\usepackage{lmodern}
\usepackage{microtype}
\usepackage[a4paper,margin=27mm,headheight=14pt]{geometry}
\usepackage{amsmath,amssymb,amsthm,mathtools}
\usepackage{booktabs,longtable,array}
\usepackage{enumitem}
\usepackage{xcolor}
\usepackage{pid-rs-report-tables}
\usepackage{hyperref}
\usepackage{fancyhdr}

\setlength{\emergencystretch}{2em}
\newcolumntype{L}[1]{>{\raggedright\arraybackslash}p{#1}}

\ifdefined\pdfinfoomitdate
  \pdfinfoomitdate=1
\fi
\ifdefined\pdftrailerid
  \pdftrailerid{}
\fi
\ifdefined\pdfsuppressptexinfo
  \pdfsuppressptexinfo=15
\fi

\hypersetup{%
  colorlinks=true,
  linkcolor=blue!55!black,
  citecolor=blue!55!black,
  urlcolor=blue!55!black,
  pdftitle={Finite-Alphabet Plug-In Convergence for Shared-Exclusions PID},
  pdfsubject={Exact-real theorem, quantitative bounds, and evidence boundary},
  pdfkeywords={partial information decomposition, shared exclusions, finite alphabet, convergence},
  pdfauthor={pid-rs project analysis}
}

\PidConfigureSections
\PidSetRunningHeads{pid-rs finite-alphabet convergence}{project validation}
\PidConfigureAbstract
\PidConfigureLinksAndListings

\newtheorem{theorem}{Theorem}[section]
\newtheorem{corollary}[theorem]{Corollary}
\newtheorem{lemma}[theorem]{Lemma}
\newtheorem{proposition}[theorem]{Proposition}
\theoremstyle{definition}
\newtheorem{definition}[theorem]{Definition}
\newtheorem{remark}[theorem]{Remark}

\newcommand{\cZ}{\mathcal{Z}}
\newcommand{\cS}{\mathcal{S}}
\newcommand{\cT}{\mathcal{T}}
\newcommand{\supp}{\operatorname{supp}}
\newcommand{\Phat}{\widehat{P}}
\newcommand{\E}{\mathbb{E}}
\newcommand{\Prb}{\mathbb{P}}
\newcommand{\R}{\mathbb{R}}
\newcommand{\Ind}{\mathbf{1}}
\newcommand{\normone}[1]{\left\lVert #1\right\rVert_1}
\newcommand{\norminf}[1]{\left\lVert #1\right\rVert_\infty}
\newcommand{\sx}{\mathrm{sx}}
\newcommand{\imin}{I_{\min}}

\title{\PidReportTitle
  {Finite-Alphabet Plug-In Convergence\\for Shared-Exclusions PID}
  {Exact-real theorem, local moduli, an anytime i.i.d.\ envelope,\\
   and a formal evidence boundary}}
\author{pid-rs project analysis}
\date{22 July 2026}

\begin{document}
\PidMakeTitle

\begin{abstract}
This document gives new project-defined validation for existing information measures in
\texttt{pid-rs}.  Its primary object is the Wibral-lineage categorical shared-exclusions partial
information decomposition (SxPID).  The paper-defined functional is not new.  The document proves
exact-real plug-in convergence on fixed finite alphabets, gives explicit local continuity moduli,
and composes them with a conservative time-uniform i.i.d.\ concentration event.  A secondary
corollary covers the Williams--Beer $\imin$ functional and fixed Shannon combinations.  A
frozen-transform corollary states explicit sufficient conditions under which an independently trained finite
map preserves the argument.  The theorem does not prove convergence of binary64 program outputs,
dependent-window validity, drift calibration, or scientific novelty.  Invalid stronger claims and
counterexamples remain in the document.
\end{abstract}

A separate repository paper on SxPID under a dependency coloring gives a categorical result for a declared
deterministic coloring.  Its common-law theorem requires all complete rows to share one finite law.
Complete rows in each color must be mutually independent.  It also separates nonidentical-law
drift.  That result does not change the i.i.d.\ and ergodic claims in this paper, and it does not
cover generic mixing, adaptive colors, or continuous estimators.

\begingroup
\small
\setlength{\parskip}{0pt}
\tableofcontents
\endgroup

\section{Status and provenance}

The mathematical object and the validation artifact have different origins.

\begin{center}
\begin{tabular}{L{0.27\textwidth}L{0.27\textwidth}L{0.36\textwidth}}
\toprule
\PidTableHeaderRow
Object & Origin & Result in this document \\
\midrule
Categorical SxPID & Paper-defined by Makkeh, Gutknecht, and Wibral; its part-whole lattice has a
formal-logic basis due to Gutknecht, Wibral, and Makkeh~\cite{makkeh2021,gutknecht2021} & Exact-real convergence for two to four
sources, local moduli, and bounded code evidence \\
$\imin$ & Paper-defined by Williams and Beer~\cite{williamsbeer2010} & Secondary finite-minimum corollary for two and three
sources \\
Shannon entropy and mutual information & Paper-defined classical quantities~\cite{shannon1948} & Exact-real
convergence and entropy continuity bounds \\
Discrete co-information and O-information & Paper-defined quantities~\cite{bell2003,mcgill1954,rosas2019} & Convergence as fixed finite
entropy combinations \\
Target-free normalized ratios & Project-defined analogues in \texttt{pid-rs} & Positive-denominator
ratio convergence \\
Target-conditioned normalized ratios & Paper-defined ratios with a project-defined report
policy~\cite{gutknecht2025} & Positive-denominator ratio convergence and a status boundary \\
Fitted quantizer plus categorical PID & Project-defined composition; quantization has published
background~\cite{grayneuhoff1998} & Conditional frozen-transform corollary \\
\bottomrule
\end{tabular}
\end{center}

The machine-readable authority is \texttt{method-catalog.json}.  The rendered authority is
\texttt{METHODS.md}.  This document defines no new estimator or public API\@.  The exact Rust routes
covered by the result are:
\begingroup
\raggedright
\begin{itemize}[leftmargin=*,nosep]
  \item SxPID namespace \nolinkurl{stable::categorical}: \nolinkurl{discrete_sxpid2},
    \nolinkurl{discrete_sxpid2_averaged}, \nolinkurl{discrete_sxpid3},
    \nolinkurl{discrete_sxpid3_averaged}, \nolinkurl{discrete_sxpid_n}, and
    \nolinkurl{discrete_sxpid_n_averaged};
  \item $\imin$ namespace \nolinkurl{stable::imin}: \nolinkurl{imin_pid2} and
    \nolinkurl{imin_pid3};
  \item Shannon namespace \nolinkurl{diagnostics}: \nolinkurl{entropy_discrete},
    \nolinkurl{joint_entropy_discrete}, \nolinkurl{co_information_pairwise_discrete}, and
    \nolinkurl{o_information_discrete};
  \item normalized-report namespace \nolinkurl{diagnostics}: \nolinkurl{red_degree_discrete},
    \nolinkurl{vul_degree_discrete}, \nolinkurl{average_degree_of_redundancy}, and
    \nolinkurl{average_degree_of_vulnerability};
  \item fitted SxPID namespace \nolinkurl{stable::quantized}:
    \nolinkurl{fitted_quantized_sxpid2}, \nolinkurl{fitted_quantized_sxpid3}, and
    \nolinkurl{fitted_quantized_sxpid_n};
  \item fitted $\imin$ namespace \nolinkurl{stable::imin}: \nolinkurl{imin_pid2_quantized} and
    \nolinkurl{imin_pid3_quantized}.
\end{itemize}
\endgroup

\section*{Reader's guide, foundational notation, and prerequisites}
\addcontentsline{toc}{section}{Reader's guide, notation, and prerequisites}

No prior knowledge of PID is required for the arguments below.  A \emph{probability law} $Q$ on a
finite set $\cZ$ is a vector $(Q(z))_{z\in\cZ}$ with nonnegative entries summing to one.  Its
support is
\[
  \supp(Q)=\{z\in\cZ:Q(z)>0\}.
\]
For two such laws,
\[
  \normone{Q-P}=\sum_{z\in\cZ}|Q(z)-P(z)|,
  \qquad
  \operatorname{TV}(Q,P)=\frac12\normone{Q-P}
  =\sup_{E\subseteq\cZ}|Q(E)-P(E)|.
\]
Marginalization means summing joint-cell masses over coordinates that are not retained; it cannot
increase $L^1$ distance.  A \emph{plug-in} quantity is the population formula evaluated after
replacing the unknown law by an empirical law.

PID seeks coordinates whose sum reconstructs a mutual-information quantity while distinguishing
information shared by several sources, information unique to one source, and information available
only jointly.  SxPID constructs these coordinates from logical source events.  Put
$[m]=\{1,\ldots,m\}$.  A nonempty set $a\subseteq[m]$ is a \emph{source collection};
$S_a=(S_i)_{i\in a}$ and $s_a=(s_i)_{i\in a}$.  Thus
\[
  \{S_a=s_a\}=\bigcap_{i\in a}\{S_i=s_i\}.
\]
A lattice node $\alpha$ is a nonempty \emph{antichain} of nonempty source collections: no two
distinct members of $\alpha$ contain one another.  The redundancy order used here is
\[
  \alpha\preceq\beta
  \quad\Longleftrightarrow\quad
  \text{for every }b\in\beta\text{ there is an }a\in\alpha\text{ with }a\subseteq b.
\]
For $u\in\{+,-,\sx\}$, cumulative coordinates and pointwise atoms are related by the finite
poset zeta transform
\begin{equation}\label{eq:zeta-primer-finite}
  c^u_\alpha(z;Q)
  =\sum_{\beta\preceq\alpha}\pi^u_\beta(z;Q),
  \qquad
  \pi^u_\alpha(z;Q)
  =\sum_\beta M_{\alpha\beta}c^u_\beta(z;Q).
\end{equation}
Here $M$ is the unique M{\"o}bius-inversion matrix of this fixed finite order.  The averaged atom is
\begin{equation}\label{eq:average-primer-finite}
  \bar\pi^u_\alpha(Q)
  =\sum_{z\in\supp(Q)}Q(z)\pi^u_\alpha(z;Q).
\end{equation}
Equations~\eqref{eq:zeta-primer-finite}--\eqref{eq:average-primer-finite} specify what ``M{\"o}bius
atom'' and ``averaged atom'' mean everywhere below; no unstated vector-index convention is used.
The signed-net coordinate is the informative coordinate minus the misinformative coordinate, both
before and after inversion.

The proofs treat the following standard results as cited prerequisites rather than reproving them:
the strong law of large numbers, the pointwise ergodic theorem, elementary continuity and the mean
value theorem for $\log$, nonnegativity of Kullback--Leibler divergence, finite-poset M{\"o}bius
inversion, Hoeffding's inequality, and the Fannes--Audenaert entropy-continuity inequality.
Every application records its required finite-alphabet, support, independence, or denominator
premise.  ``Exact real'' means arithmetic in the mathematical real numbers; it is deliberately
distinct from a theorem about a particular binary64 implementation.

\section{Finite setting and shared-exclusions quantities}

Fix a source count $m$, finite source alphabets $\cS_1,\ldots,\cS_m$, a finite target alphabet
$\cT$, and the joint alphabet
\begin{equation}
  \cZ=\cS_1\times\cdots\times\cS_m\times\cT.
\end{equation}
Let $P$ be one probability law on $\cZ$, and put
\begin{equation}
  S=\supp(P),
  \qquad
  p_{\min}=\min_{z\in S}P(z)>0.
\end{equation}
All logarithms are natural, so all information values use nats.  The redundancy lattice is fixed
and finite.  Its nodes are nonempty antichains of nonempty source subsets.  In \texttt{pid-rs},
$m\in\{2,3,4\}$ for SxPID and $m\in\{2,3\}$ for $\imin$.

For observations $Z_1,Z_2,\ldots$, define the prefix empirical law
\begin{equation}
  \Phat_n(z)=\frac{1}{n}\sum_{j=1}^{n}\Ind\{Z_j=z\}.
\end{equation}

Fix a supported realization $z=(s_1,\ldots,s_m,t)$ and a lattice node $\alpha$.  Define
\begin{align}
  A_\alpha(s)
    &=\bigcup_{a\in\alpha}\{S_a=s_a\}, \\
  B_\alpha(s,t)
    &=A_\alpha(s)\cap\{T=t\}, \\
  C(t)
    &=\{T=t\}.
\end{align}
Each of these three events contains $z$.  Its $P$-mass is therefore at least
$P(z)\ge p_{\min}$.  For any law $Q$
whose required masses are positive, the cumulative informative and misinformative terms are
\begin{align}
  c^+_\alpha(z;Q)
    &=-\log Q(A_\alpha(s)), \\
  c^-_\alpha(z;Q)
    &=\log\frac{Q(C(t))}{Q(B_\alpha(s,t))}.
\end{align}
The cumulative net term is $c^\sx_\alpha=c^+_\alpha-c^-_\alpha$.  A fixed M{\"o}bius matrix $M$ maps
each cumulative vector to its pointwise atom vector.  The averaged atom under $Q$ is the
$Q$-weighted sum over joint realizations.

\paragraph{The two-source dictionary and a worked empirical example.}
For $m=2$, the four antichains are
\[
  \alpha_{\mathrm{red}}=\{\{1\},\{2\}\},\qquad
  \alpha_1=\{\{1\}\},\qquad
  \alpha_2=\{\{2\}\},\qquad
  \alpha_{\mathrm{syn}}=\{\{1,2\}\}.
\]
The only covering relations are
\[
  \alpha_{\mathrm{red}}\prec\alpha_1\prec\alpha_{\mathrm{syn}},
  \qquad
  \alpha_{\mathrm{red}}\prec\alpha_2\prec\alpha_{\mathrm{syn}},
\]
and $\alpha_1,\alpha_2$ are incomparable.  For the signed-net atoms, these nodes are called
redundancy, unique information from $S_1$, unique information from $S_2$, and synergy,
respectively.  For each $u\in\{+,-,\sx\}$, the zeta relation is explicitly
\begin{align}
  c^u_{\alpha_{\mathrm{red}}}
    &=\pi^u_{\alpha_{\mathrm{red}}}, \\
  c^u_{\alpha_1}
    &=\pi^u_{\alpha_{\mathrm{red}}}+\pi^u_{\alpha_1}, \\
  c^u_{\alpha_2}
    &=\pi^u_{\alpha_{\mathrm{red}}}+\pi^u_{\alpha_2}, \\
  c^u_{\alpha_{\mathrm{syn}}}
    &=\pi^u_{\alpha_{\mathrm{red}}}+\pi^u_{\alpha_1}
      +\pi^u_{\alpha_2}+\pi^u_{\alpha_{\mathrm{syn}}}.
\end{align}

As a complete empirical calculation, take the four binary XOR rows
\[
  (s_1,s_2,t)\in\{(0,0,0),(0,1,1),(1,0,1),(1,1,0)\},
\]
each once, and write $Q=\widehat P_4$.  Thus
\[
  Q(s_1,s_2,t)
  =\begin{cases}
      1/4,&t=s_1\mathbin{\oplus}s_2,\\
      0,&\text{otherwise}.
    \end{cases}
\]
For every supported row, $Q(C(t))=1/2$, and the cumulative values are
\begin{center}
\small
\begin{tabular}{lccccc}
\toprule
\PidTableHeaderRow
Node & $Q(A_\alpha)$ & $Q(B_\alpha)$ & $c^+_\alpha$ & $c^-_\alpha$
  & $c^\sx_\alpha$ \\
\midrule
$\alpha_{\mathrm{red}}$ & $3/4$ & $1/4$ & $\log(4/3)$ & $\log2$ & $\log(2/3)$ \\
$\alpha_1$              & $1/2$ & $1/4$ & $\log2$     & $\log2$ & $0$ \\
$\alpha_2$              & $1/2$ & $1/4$ & $\log2$     & $\log2$ & $0$ \\
$\alpha_{\mathrm{syn}}$ & $1/4$ & $1/4$ & $\log4$     & $\log2$ & $\log2$ \\
\bottomrule
\end{tabular}
\end{center}
M{\"o}bius inversion therefore gives
\begin{center}
\small
\begin{tabular}{lccc}
\toprule
\PidTableHeaderRow
Atom & $\pi^+$ & $\pi^-$ & $\pi^\sx=\pi^+-\pi^-$ \\
\midrule
Redundancy     & $\log(4/3)$ & $\log2$ & $\log(2/3)$ \\
Unique $S_1$   & $\log(3/2)$ & $0$     & $\log(3/2)$ \\
Unique $S_2$   & $\log(3/2)$ & $0$     & $\log(3/2)$ \\
Synergy        & $\log(4/3)$ & $0$     & $\log(4/3)$ \\
\bottomrule
\end{tabular}
\end{center}
All four supported rows give the same pointwise values, so these are also the averaged atoms.
They reconstruct
\[
  \log(2/3)+2\log(3/2)+\log(4/3)
  =\log2
  =I_Q((S_1,S_2);T).
\]
In particular, the negative redundancy is not clipped: together with either positive unique atom
it gives $I_Q(S_i;T)=0$.

\paragraph{Williams--Beer \texorpdfstring{$\imin$}{I-min} coordinates.}
For a nonempty source collection $a$ and a target value with $Q(t)>0$, define the
source-specific information
\begin{equation}\label{eq:ispec}
  J_a(t;Q)
  :=D\!\left(Q_{S_a\mid T=t}\,\middle\|\,Q_{S_a}\right)
  =\sum_{s_a:Q(s_a,t)>0}
    \frac{Q(s_a,t)}{Q(t)}
    \log\frac{Q(s_a,t)}{Q(s_a)Q(t)}.
\end{equation}
For an antichain $\alpha$, the target-specific minimum and its target average are
\begin{align}
  R_\alpha(t;Q)
    &:=\min_{a\in\alpha}J_a(t;Q), \\
  I^{\min}_{\cap}(\alpha;Q)
    &:=\sum_{t:Q(t)>0}Q(t)R_\alpha(t;Q).
\end{align}
The values $I^{\min}_{\cap}(\alpha;Q)$ are the cumulative Williams--Beer redundancies on the same
finite antichain order.  Their PID atoms are
\begin{equation}
  \pi^{\min}_\alpha(Q)
  :=\sum_\beta M_{\alpha\beta}I^{\min}_{\cap}(\beta;Q),
\end{equation}
where $M$ is the M{\"o}bius-inversion matrix from
Equation~\eqref{eq:zeta-primer-finite}.  Zero-mass target values are omitted.  A finite minimum of
continuous coordinates remains continuous at a tie, although it need not be differentiable.

\section{Deterministic convergence and sampling corollary}

\begin{theorem}[Deterministic plug-in implication]\label{thm:deterministic}
Let $Q_n$ be probability laws on $\cZ$.  Assume
\begin{enumerate}[label=(\roman*)]
  \item $Q_n(z)\to P(z)$ for every $z\in\cZ$; and
  \item $\supp(Q_n)\subseteq S$ for every $n$.
\end{enumerate}
Then there is a finite $N_0$ such that $\supp(Q_n)=S$ for every $n\ge N_0$.  On this tail, all
supported-event denominators are positive.  The following exact-real quantities converge to their
corresponding values under $P$:
\begin{enumerate}
  \item every supported-realization SxPID cumulative informative and misinformative term;
  \item every informative, misinformative, and signed-net M{\"o}bius atom;
  \item every probability-weighted averaged SxPID atom;
  \item every supported-target $\imin$ specific-information value and every fixed two- or
    three-source $\imin$ atom; and
  \item every fixed finite-alphabet Shannon entropy, mutual information, co-information, and
    O-information expression formed from the same law.
\end{enumerate}
Pointwise SxPID convergence is keyed by the joint realization, not by a returned-vector index.
\end{theorem}

\begin{proof}
For each $z\in S$, coordinate convergence and $P(z)>0$ give $Q_n(z)>P(z)/2$ for all sufficiently
large $n$.  The support is finite, so one finite $N_0$ works for every supported cell.  Assumption
(ii) excludes new cells.  This proves support equality on the tail.

Fix $z$ and $\alpha$.  Event masses are finite linear sums of cell masses, so the masses of
$A_\alpha$, $B_\alpha$, and $C$ converge.  Their limits are at least $P(z)>0$.  Continuity of the
logarithm on $(0,\infty)$ gives convergence of $c^+_\alpha$ and $c^-_\alpha$.  The fixed finite
	linear map $M$ preserves convergence.  Subtraction gives the net atoms.  After support
	stabilization, an average is a fixed finite sum whose weights and pointwise values both converge.
	No off-support $x\log x$ argument is needed for this SxPID step because Assumption~(ii) and
	eventual support equality keep the average on one fixed finite support face.

For the specific-information values in Equation~\eqref{eq:ispec}, support stabilization fixes
every positive marginal cell and makes every required denominator positive.  Each finite term is
continuous.  A finite minimum of continuous functions is continuous, including at a tie.  Target
averaging and the fixed M{\"o}bius map preserve convergence.  A tie can remove differentiability,
but it does not remove continuity.

Finally, define $h(0)=0$ and $h(x)=-x\log x$ for $x>0$.  The function $h$ is continuous on
$[0,1]$.  Each finite-alphabet entropy is a finite sum of $h$.  Marginalization and fixed finite
linear combinations preserve convergence.
\end{proof}

\begin{corollary}[I.i.d.\ and stationary ergodic sampling]\label{cor:sampling}
The assumptions of Theorem~\ref{thm:deterministic} hold almost surely when either
\begin{enumerate}[label=(\roman*)]
  \item the observations are i.i.d.\ with law $P$; or
  \item the process is strictly stationary and ergodic with one-time marginal law $P$.
\end{enumerate}
Strict stationarity without ergodicity is not sufficient.
\end{corollary}

\begin{proof}
Apply the strong law, or the pointwise ergodic theorem, to the indicator of each cell.  A finite
intersection gives simultaneous coordinate convergence.  If $P(z)=0$, every time index has zero
probability of that cell.  A countable union shows that no such cell occurs, almost surely.  Thus,
the empirical support is always a subset of $S$ on one probability-one event.
\end{proof}

\section{Quantitative local moduli for SxPID}

Let $p$ be the probability vector of $P$, and let $q$ be another probability law on the same finite
joint alphabet.  Define
\begin{equation}
  \delta=\normone{q-p},
  \qquad
  L=\log(1/p_{\min}),
\end{equation}
and assume $\delta\le p_{\min}/2$.  For every event $E$,
\begin{equation}\label{eq:event-loose}
  |q(E)-p(E)|\le\delta.
\end{equation}
This deliberately loose inequality is sufficient for the following stated constants.

\begin{theorem}[SxPID local bounds]\label{thm:sx-local}
For each supported realization and lattice node,
\begin{align}
  |c^+_\alpha(z;q)-c^+_\alpha(z;p)|
    &\le \frac{2\delta}{p_{\min}}, \\
  |c^-_\alpha(z;q)-c^-_\alpha(z;p)|
    &\le \frac{4\delta}{p_{\min}}.
\end{align}
Let
\begin{equation}
  \norminf{M}=\max_i\sum_j|M_{ij}|.
\end{equation}
The pointwise atom errors obey
\begin{align}\label{eq:point-bounds}
  \text{informative:}\quad
  |\Delta\pi^+|&\le\frac{2\norminf{M}\delta}{p_{\min}}, \\
  \text{misinformative:}\quad
  |\Delta\pi^-|&\le\frac{4\norminf{M}\delta}{p_{\min}}, \\
  \text{net:}\quad
  |\Delta\pi^\sx|&\le\frac{6\norminf{M}\delta}{p_{\min}}.
\end{align}

If also $\supp(q)\subseteq S$, the averaged errors obey
\begin{align}
  \text{informative:}\quad
  |\Delta\bar\pi^+|&\le\norminf{M}\left(\frac{2}{p_{\min}}+L\right)\delta, \\
  \text{misinformative:}\quad
  |\Delta\bar\pi^-|&\le\norminf{M}\left(\frac{4}{p_{\min}}+L\right)\delta, \\
  \text{net:}\quad
  |\Delta\bar\pi^\sx|&\le\norminf{M}\left(\frac{6}{p_{\min}}+L\right)\delta.
\end{align}
\end{theorem}

\begin{proof}
Each required event $E$ has $p(E)\ge p_{\min}$.  Equations~\eqref{eq:event-loose} and
$\delta\le p_{\min}/2$ give $q(E)\ge p_{\min}/2$.  The mean-value theorem for $\log$ gives
\begin{equation}
  |\log q(E)-\log p(E)|
  \le \frac{|q(E)-p(E)|}{\min\{q(E),p(E)\}}
  \le \frac{2\delta}{p_{\min}}.
\end{equation}
The informative term contains one logarithm.  The misinformative term contains two.  The row-sum
norm of $M$ gives the first two atom bounds in~\eqref{eq:point-bounds}; the triangle inequality
gives the net bound.

For one atom component $\pi$, decompose its averaged error as
\begin{align}
  \left|\sum_{z\in S}q(z)\pi(z;q)-\sum_{z\in S}p(z)\pi(z;p)\right|
  &\le \sum_{z\in S}q(z)|\pi(z;q)-\pi(z;p)| \\
  &\quad +\sum_{z\in S}|q(z)-p(z)|\,|\pi(z;p)|.
\end{align}
At $p$, each cumulative informative and misinformative coordinate lies in $[0,L]$.  For the net
term, $c^+_\alpha-c^-_\alpha\in[-L,L]$.  Apply the row-sum norm once.  This gives one factor of $L$
for each of the three weight-change terms and proves the averaged bounds.
\end{proof}

\begin{remark}[Why the common-support condition matters]
Under $\delta\le p_{\min}/2$, the condition $\supp(q)\subseteq S$ implies $\supp(q)=S$.  Without
common support, a concrete two-source law shows the obstruction.  Let all three alphabets be
binary, let
\[
  P=\delta_{(0,0,0)},
  \qquad
  q_\varepsilon
    =(1-\varepsilon)\delta_{(0,0,0)}
     +\varepsilon\delta_{(1,1,1)},
  \qquad 0<\varepsilon<\frac12.
\]
Then
\[
  \operatorname{TV}(q_\varepsilon,P)=\varepsilon,
  \qquad
  \normone{q_\varepsilon-P}=2\varepsilon.
\]
At the bottom node $\alpha_{\mathrm{red}}=\{\{1\},\{2\}\}$, the shared source event for
$(0,0,0)$ has mass $1-\varepsilon$, while the event for $(1,1,1)$ has mass $\varepsilon$.
The misinformative term is zero at both rows, and the averaged informative---hence also net---bottom
atom is
\[
  \bar\pi^\sx_{\alpha_{\mathrm{red}}}(q_\varepsilon)
  =-(1-\varepsilon)\log(1-\varepsilon)-\varepsilon\log\varepsilon
  =h_2(\varepsilon),
\]
whereas its value under $P$ is zero.  Therefore
\[
  \frac{
    \left|\bar\pi^\sx_{\alpha_{\mathrm{red}}}(q_\varepsilon)
          -\bar\pi^\sx_{\alpha_{\mathrm{red}}}(P)\right|
  }{\normone{q_\varepsilon-P}}
  =\frac{h_2(\varepsilon)}{2\varepsilon}
  \longrightarrow\infty.
\]
Thus no support-independent linear constant can replace the common-support premise.
\end{remark}

These are exact-real law-level inequalities.  They do not include binary64 rounding, platform
logarithm variation, resource rejection, or implementation refinement.  The Rust SxPID and $\imin$
paths reject sample counts above $2^{53}$.

\section{Secondary \texorpdfstring{$\imin$}{I-min} modulus}

This section is a secondary validation control.  It is not the research priority of the document.
Assume $\supp(q)\subseteq S$ and $\delta\le p_{\min}/2$.  For a supported target $t$, let
\begin{equation}
  r=q(s_a,t),\qquad u=q(s_a),\qquad v=q(t),\qquad
  w=\frac{r}{v},\qquad \ell=\log\frac{r}{uv}.
\end{equation}
Marginalization contracts joint $L^1$ distance.  Every positive $p$-marginal is at least
$p_{\min}$, and every corresponding $q$-marginal is at least $p_{\min}/2$.  Hence
\begin{equation}\label{eq:weight-bound}
  \sum_{s_a}|w_q(s_a)-w_p(s_a)|
  \le \frac{4\delta}{p_{\min}},
\end{equation}
and
\begin{equation}\label{eq:log-bound}
  |\ell_q-\ell_p|\le\frac{6\delta}{p_{\min}}.
\end{equation}
To see~\eqref{eq:weight-bound}, put $r_q(s_a)=q(s_a,t)$, $v_q=q(t)$, and write
\begin{align}
  \sum_{s_a}\left|\frac{r_q(s_a)}{v_q}-\frac{r_p(s_a)}{v_p}\right|
  &\le
  \frac{\sum_{s_a}|r_q(s_a)-r_p(s_a)|}{v_q}
  +\frac{|v_q-v_p|}{v_q} \\
  &\le \frac{2\delta}{p_{\min}/2}.
\end{align}
	Each of the three logarithm factors changes by at most $2\delta/p_{\min}$.
	Equation~\eqref{eq:log-bound} is their sum.  Also,
	$r_p,u_p,v_p\in[p_{\min},1]$ whenever $r_p>0$, and $r_p\le u_p,v_p$.  Hence
	$p_{\min}\le r_p/(u_pv_p)\le1/p_{\min}$ and
	$|\ell_p|\le L\le\log(2/p_{\min})$.  Splitting the specific-information difference gives
	\begin{align}
	  |J_a(t;q)-J_a(t;p)|
	  &\le
	  \sum_{s_a}w_q(s_a)|\ell_q(s_a)-\ell_p(s_a)|\\
	  &\quad+
	  \sum_{s_a}|w_q(s_a)-w_p(s_a)|\,|\ell_p(s_a)|.
	\end{align}
	The first sum is at most $6\delta/p_{\min}$ because $w_q$ is a probability vector.  The second
	is at most $(4\delta/p_{\min})\log(2/p_{\min})$ by
	Equation~\eqref{eq:weight-bound}.  Therefore,
\begin{equation}\label{eq:imin-bound}
  |J_a(t;q)-J_a(t;p)|
  \le
  \left[4\log\left(\frac{2}{p_{\min}}\right)+6\right]
  \frac{\delta}{p_{\min}}.
\end{equation}
A finite minimum is one-Lipschitz in the sup norm.  If $C_J$ is the coefficient of $\delta$ in
\eqref{eq:imin-bound}, define
\begin{align}
  R_\alpha(t;Q)&=\min_{a\in\alpha}J_a(t;Q), \\
  I^{\min}_{\cap}(\alpha;Q)&=\sum_t Q(t)R_\alpha(t;Q).
\end{align}
Each $J_a(t;p)$ is a Kullback--Leibler divergence, so it is nonnegative.  Also,
\begin{equation}
  J_a(t;p)
  =\sum_{s_a}p(s_a\mid t)\log\frac{p(t\mid s_a)}{p(t)}
  \le -\log p(t)
  \le L.
\end{equation}
Thus, $0\le R_\alpha(t;p)\le L$.  Marginalization gives
$\normone{q_T-p_T}\le\delta$.  Therefore,
\begin{align}
  \left|I^{\min}_{\cap}(\alpha;q)-I^{\min}_{\cap}(\alpha;p)\right|
  &\le \sum_t q(t)|R_\alpha(t;q)-R_\alpha(t;p)| \\
  &\quad +\sum_t|q(t)-p(t)|R_\alpha(t;p) \\
  &\le (C_J+L)\delta.
\end{align}
Since $L=\log(1/p_{\min})$, the last line is the target-averaged redundancy bound.
The fixed $\imin$ M{\"o}bius matrix multiplies this bound by its absolute row-sum norm.  A minimum tie
preserves continuity but can remove differentiability and a generic normal limit.

\paragraph{A complete minimum-tie crossing.}
Let $T$ be a uniform bit.  For $|\theta|<1/4$, let
\[
  e_1(\theta)=\frac14+\theta,
  \qquad
  e_2(\theta)=\frac14-\theta,
\]
let $N_1\sim\operatorname{Bernoulli}(e_1(\theta))$ and
$N_2\sim\operatorname{Bernoulli}(e_2(\theta))$ be mutually independent and independent of $T$,
and put
\[
  S_i=T\mathbin{\oplus}N_i.
\]
Equivalently, with $b_e(0)=1-e$ and $b_e(1)=e$, the full-support joint law is
\[
  P_\theta(s_1,s_2,t)
  =\frac12\,b_{e_1(\theta)}(s_1\mathbin{\oplus}t)\,
             b_{e_2(\theta)}(s_2\mathbin{\oplus}t).
\]
Each $S_i$ is uniform, and for either target value
\[
  J_{\{i\}}(t;P_\theta)
  =f(e_i(\theta)),
  \qquad
  f(e):=\log2-h_2(e),
\]
where $h_2(e)=-e\log e-(1-e)\log(1-e)$.  Since
\[
  f'(e)=\log\frac{e}{1-e}<0
  \quad\text{for }0<e<\frac12,
\]
the redundancy node $\alpha_{\mathrm{red}}=\{\{1\},\{2\}\}$ satisfies
\[
  R_{\alpha_{\mathrm{red}}}(t;P_\theta)
  =\min\{f(e_1(\theta)),f(e_2(\theta))\}
  =f\!\left(\frac14+|\theta|\right).
\]
The source-specific values tie at $\theta=0$ and exchange minimizers across zero.  Since
$f'(1/4)=-\log3$, the left derivative is $+\log3$ and the right derivative is $-\log3$.
Target averaging does not change this expression, so the averaged $\imin$ redundancy is
continuous but nondifferentiable at the crossing.

\section{Shannon and ratio bounds}

\paragraph{Definitions and sign conventions.}
For a finite tuple $X_A$, write
\[
  H_Q(X_A):=-\sum_{x_A:Q(x_A)>0}Q(x_A)\log Q(x_A),
  \qquad H_Q(\varnothing)=0.
\]
Mutual information and conditional mutual information are the entropy combinations
\begin{align}
  I_Q(X;Y)
    &=H_Q(X)+H_Q(Y)-H_Q(X,Y),\\
  I_Q(X;Y\mid Z)
    &=H_Q(X,Z)+H_Q(Y,Z)-H_Q(Z)-H_Q(X,Y,Z).
\end{align}
The implementation uses Bell's co-information sign:
\begin{align}
  \operatorname{CoI}_{\mathrm{Bell}}(X_1,X_2;T)
    &:=I(X_1;T)+I(X_2;T)-I((X_1,X_2);T)\\
    &=I(X_1;X_2)-I(X_1;X_2\mid T)\\
    &=H(X_1)+H(X_2)+H(T)-H(X_1,X_2)-H(X_1,T)\notag\\
    &\qquad{}-H(X_2,T)+H(X_1,X_2,T).
\end{align}
The McGill-style interaction-information convention recorded by this project has the opposite
sign:
\[
  \operatorname{II}_{\mathrm{McGill}}
  =-\operatorname{CoI}_{\mathrm{Bell}}.
\]
Thus the Bell-sign co-information of two independent uniform bits and their XOR target is
$-\log2$.

For $r\ge2$, with $X_{-i}$ denoting all variables except $X_i$, the O-information is
\[
  \Omega(X_1,\ldots,X_r)
  :=(r-2)H(X_1,\ldots,X_r)
    +\sum_{i=1}^{r}H(X_i)
    -\sum_{i=1}^{r}H(X_{-i}).
\]
For $r=3$, this equals the Bell-sign co-information.  Positive O-information is conventionally
called redundancy-dominated and negative O-information synergy-dominated; it is not itself a PID
atom.

The project-defined target-free ratios are
\begin{align}
  \mathrm{Red}^{\circ}(X_1,\ldots,X_r)
    &:=\frac{\sum_iH(X_i)}{H(X_1,\ldots,X_r)},\\
  \mathrm{Vul}^{\circ}(X_1,\ldots,X_r)
    &:=\frac{\sum_iH(X_i\mid X_{-i})}{H(X_1,\ldots,X_r)}\\
    &=\frac{rH(X_1,\ldots,X_r)-\sum_iH(X_{-i})}
            {H(X_1,\ldots,X_r)}.
\end{align}
They require positive joint entropy.  The published target-conditioned ratios are
\begin{align}
  \bar r(T;S_1,\ldots,S_m)
    &:=\frac{\sum_iI(T;S_i)}{I(T;S_1,\ldots,S_m)},\\
  \bar v(T;S_1,\ldots,S_m)
    &:=\frac{\sum_iI(T;S_i\mid S_{-i})}{I(T;S_1,\ldots,S_m)}\\
    &=\frac{\sum_i\left[I(T;S_1,\ldots,S_m)-I(T;S_{-i})\right]}
            {I(T;S_1,\ldots,S_m)}.
\end{align}
Here $I(T;\varnothing)=0$.  These ratios require positive joint mutual information.  All four
ratios are unitless.  Any reporting threshold is a software policy, not part of the mathematical
definitions.

Let one entropy term have alphabet size $K\ge2$, and put
\begin{equation}
  \varepsilon=\tfrac12\normone{q-p}.
\end{equation}
Apply Audenaert's bound~\cite{audenaert2007} to diagonal density matrices with diagonals $p$ and
$q$.  Their trace distance is $\varepsilon$, and their von Neumann entropies are $H(p)$ and $H(q)$.
For $0\le\varepsilon\le1-1/K$, the Fannes--Audenaert bound gives
\begin{equation}\label{eq:fannes}
  |H(q)-H(p)|
  \le g_K(\varepsilon)
  :=\varepsilon\log(K-1)+h_2(\varepsilon),
\end{equation}
where
\begin{equation}
  h_2(x)=-x\log x-(1-x)\log(1-x),
  \qquad h_2(0)=h_2(1)=0.
\end{equation}
	The function $g_K$ is increasing only up to $1-1/K$.  For a marginal $A$ with alphabet size
	$K_A\ge2$ and only a joint variation bound available, the safe bound is
\begin{equation}\label{eq:fannes-marginal}
  |H(q_A)-H(p_A)|
  \le
  g_{K_A}\!\left(\min\{\varepsilon,1-1/K_A\}\right).
	\end{equation}
	If the marginal variation $\varepsilon_A$ itself is available and
	$\varepsilon_A\le1-1/K_A$, use $g_{K_A}(\varepsilon_A)$.  If
	$\varepsilon_A>1-1/K_A$, use the valid cap $\log K_A$.  If $K_A=1$, the entropy difference is
		zero.  Marginalization gives $\varepsilon_A\le\varepsilon$, which is why
		Equation~\eqref{eq:fannes-marginal} is safe.  A fixed entropy linear combination is bounded by
		the sum of the absolute coefficients times the term-specific bounds.

\paragraph{Why clipping is necessary.}
Direct substitution of the joint variation can fail.  On the four cells
$(A,B)=(0,0),(0,1),(1,0),(1,1)$, take
\[
  p=(1,0,0,0),
  \qquad
  q=\left(\frac3{10},\frac15,\frac12,0\right).
\]
Then
\[
  \operatorname{TV}(p,q)
  =\frac12\left(\frac7{10}+\frac15+\frac12\right)
  =\frac7{10},
\]
whereas
\[
  p_A=(1,0),\qquad q_A=(1/2,1/2),\qquad
  \operatorname{TV}(p_A,q_A)=\frac12.
\]
Consequently,
\[
  |H(q_A)-H(p_A)|=\log2.
\]
For a binary alphabet, $g_2(x)=h_2(x)$, and direct un-clipped substitution would give
\[
  g_2(0.7)=h_2(0.7)\approx0.610864
  <\log2\approx0.693147,
\]
which is not an upper bound.  Clipping at $1-1/2=1/2$ instead gives
$g_2(1/2)=\log2$, exactly as required.

\begin{lemma}[Ratio perturbation]\label{lem:ratio}
Let $D_p\ge d>0$, $|D_q-D_p|\le d/2$, and $|N_p|\le B$.  Then
\begin{equation}
  \left|\frac{N_q}{D_q}-\frac{N_p}{D_p}\right|
  \le
  \frac{2|N_q-N_p|}{d}
  +\frac{2B|D_q-D_p|}{d^2}.
\end{equation}
\end{lemma}

\begin{proof}
The denominator condition gives $D_q\ge d/2$.  Decompose the difference into
$(N_q-N_p)/D_q+N_p(D_p-D_q)/(D_qD_p)$ and apply the triangle inequality.
\end{proof}

Target-free normalized limits require positive population joint entropy.  Target-conditioned
normalized limits require positive population joint mutual information.  Every target-conditioned
input term must describe the same law, target, source order, units, preprocessing, evaluation
sample, and estimand.  For the convergence claim, each input sequence must be the corresponding
exact-real finite-alphabet plug-in term formed from the same empirical law.  A positive denominator
alone does not establish this coherence.

For a report threshold $\tau$, a population denominator strictly above $\tau$ guarantees eventual
\texttt{Defined} status.  A denominator below $\tau$ guarantees an eventual status other than
\texttt{Defined}.
Equality gives no eventual status guarantee.

\section{A time-uniform i.i.d.\ envelope}

This section uses Hoeffding's published inequality~\cite{hoeffding1963} and union bounds.  Their
composition here is project-defined validation and makes no scientific-novelty claim.

\begin{theorem}[All-prefix empirical-law bound]\label{thm:anytime}
Assume i.i.d.\ sampling from one fixed $K$-cell law.  Let $0<\alpha<1$ and
\begin{equation}
  \alpha_n=\frac{6\alpha}{\pi^2n^2}.
\end{equation}
With probability at least $1-\alpha$, simultaneously for every $n\ge1$,
\begin{equation}\label{eq:Bn}
  \normone{\Phat_n-P}
  \le B_n
  :=K\sqrt{\frac{\log(K\pi^2n^2/(3\alpha))}{2n}}.
\end{equation}
\end{theorem}

\begin{proof}
For one cell, Hoeffding's inequality gives
\begin{equation}
  \Prb\bigl(|\Phat_n(z)-P(z)|\ge x\bigr)\le2e^{-2nx^2}.
\end{equation}
Set
\begin{equation}
  x_n=\sqrt{\frac{\log(K\pi^2n^2/(3\alpha))}{2n}}.
\end{equation}
The failure probability for one cell at time $n$ is at most $\alpha_n/K$.  A union bound over $K$
cells and all $n$ costs
\begin{equation}
  \sum_{n\ge1}\alpha_n
  =\frac{6\alpha}{\pi^2}\sum_{n\ge1}\frac1{n^2}
  =\alpha.
\end{equation}
On the complement, summing the $K$ coordinate bounds gives~\eqref{eq:Bn}.
\end{proof}

Define
\begin{equation}
  N_0=\min\!\bigl\lbrace r\ge1:B_n\le p_{\min}/2\;\text{for every}\;n\ge r\bigr\rbrace.
\end{equation}
The integer $N_0$ is finite because $B_n\to0$.  Intersect the concentration event with the
probability-one support-containment event of Corollary~\ref{cor:sampling}.  On the intersection,
support equality holds for all $n\ge N_0$.  Substitute $B_n$ for the actual $\delta$ in
Theorem~\ref{thm:sx-local}.

	For a fixed prefix $N$, a separate support-discovery bound is
\begin{equation}
  \Prb\bigl(\supp(\Phat_N)\ne S\bigr)
  \le |S|e^{-Np_{\min}}
  \le Ke^{-Np_{\min}}.
	\end{equation}
	Indeed, one supported cell $z$ is missing after $N$ i.i.d.\ rows with probability
	$(1-P(z))^N\le e^{-NP(z)}\le e^{-Np_{\min}}$; a union bound over $S$ gives the first inequality.
	Once a state enters a cumulative prefix, it remains in every later prefix.  This persistence
	does not hold for a sliding window.
	A usable $N_0$ needs a known positive lower bound on $p_{\min}$.  An empirical minimum cannot
	replace it because an unobserved rare state can have positive population mass.  The simultaneous
	$L^1$ bound can be read at an almost-surely finite data-dependent index.  Support and local-bound
	conclusions additionally need that index to be at least $N_0$.  The law must remain the same
	fixed i.i.d.\ law.  The statement does not cover sliding or otherwise dependent windows, drift,
	feedback that changes the sampling law, rejection-selected samples, or stationary ergodic data
	without an additional rate assumption.

The theorem controls an exact-real empirical law.  It is not a confidence sequence for Rust
output.  It does not include floating-point error, the platform logarithm, resource failure,
dependence, drift, rejection selection, or the $2^{53}$ sample-count limit.

\section{Independent frozen transforms}

Let $(\Omega,\mathcal F,\Prb)$ be the underlying probability space and let
$\mathcal G\subseteq\mathcal F$ be generated by the training artifact.  Let $\mathcal X$ be
standard Borel, let $\cZ$ be finite, and let $\bot\notin\cZ$.  Suppose
\[
  \Phi:\Omega\times\mathcal X\longrightarrow\cZ\cup\{\bot\}
\]
is $\mathcal G\otimes\mathcal B(\mathcal X)$-measurable.  The same frozen map
$\Phi_\omega(x):=\Phi(\omega,x)$ is used for every evaluation row and every prefix.

Let $W_1,W_2,\ldots$ be conditionally i.i.d.\ given $\mathcal G$, with common regular conditional
law $\kappa_\omega$ on $\mathcal X$.  An i.i.d.\ raw sequence independent of the training artifact
is the special case in which $\kappa_\omega$ is a fixed law.  Put
\[
  Y_j(\omega):=\Phi(\omega,W_j(\omega))
\]
and assume
\begin{equation}\label{eq:frozen-total}
  \kappa_\omega\!\left(\{x:\Phi_\omega(x)=\bot\}\right)=0
  \quad\text{almost surely}.
\end{equation}

\begin{corollary}[Conditional push-forward target]\label{cor:transform}
Define
\[
  P_\Phi(\omega;z)
  :=\int_{\mathcal X}\Ind\{\Phi_\omega(x)=z\}\,\kappa_\omega(dx),
  \qquad z\in\cZ.
\]
On an event of probability one,
\[
  \widehat P^\Phi_n(z)
  :=\frac1n\sum_{j=1}^{n}\Ind\{Y_j=z\}
  \longrightarrow P_\Phi(z)
  \qquad\text{for every }z\in\cZ.
\]
The empirical support is always contained in $\supp(P_\Phi)$ and equals it for all sufficiently
large $n$.  Hence Theorem~\ref{thm:deterministic} applies pathwise with population law $P_\Phi$.
The limit is generally not the same functional evaluated at the unconditional mixture
$\bar P_\Phi(z):=\E[P_\Phi(z)]$.
\end{corollary}

\begin{proof}
Conditional i.i.d.\ means that, for every $k$ and $z_1,\ldots,z_k\in\cZ$,
\begin{align}
 &\Prb(Y_1=z_1,\ldots,Y_k=z_k\mid\mathcal G)(\omega)\\
 &\quad=
 \int_{\mathcal X^k}
   \prod_{j=1}^{k}\Ind\{\Phi_\omega(x_j)=z_j\}\,
   \kappa_\omega^{\otimes k}(dx_{1:k})
 =\prod_{j=1}^{k}P_\Phi(\omega;z_j).
\end{align}
Thus, on almost every conditioning fiber, the transformed sequence is i.i.d.\ with the fixed
finite law $P_\Phi(\omega;\cdot)$.  Equation~\eqref{eq:frozen-total} and a countable union ensure
that every transformed row lies in $\cZ$.

The ordinary strong law on each product-law fiber gives simultaneous convergence of the finitely
many cell indicators.  A zero-probability cell is never observed on that fiber, while convergence
to a positive cell probability makes its empirical frequency positive eventually.  This proves
support containment and stabilization.  Applying Theorem~\ref{thm:deterministic} fiberwise and
integrating the resulting conditional probability-one event proves the claim.
\end{proof}

	An \nolinkurl{OutOfRangePolicy::Error} transform satisfies the corollary only when the
	conditional evaluation law has mass one inside every inclusive fitted range and satisfies every
	input-validity rule.  If it violates the zero-failure assumption, let
	$\eta=\eta(\mathcal{G})$ be its conditional per-row failure probability.  For a training outcome
	with $\eta>0$, the first $n$ rows all succeed with probability
	$(1-\eta)^n$, and the infinite prefix fails with conditional probability one.  Unconditionally,
	this failure mechanism has probability $\Prb(\eta>0)$.  It has probability one only when
	$\eta>0$ almost surely.

Conditioning a complete prefix on success does not preserve the original-law claim.  Explicit
rejection sampling can produce a sequence that is conditionally i.i.d.\ given $\mathcal{G}$ from a
conditional law when success has positive probability.  That sequence has a different estimand and
needs a declared sampling contract.
	\nolinkurl{OutOfRangePolicy::ClampToBoundary} supplies a total range map for otherwise valid
	finite inputs, but it defines a tail-clamped categorical estimand.  It is not the error-policy
	estimand.

	Same-row fitting, changing transforms, arbitrary pooling across different fitted maps, and
	target-adaptive fitting on evaluation rows are outside the corollary.  Training-target-adaptive
	partial least squares is compatible only when fitted on the independent training artifact,
	frozen, and applied to disjoint conditionally i.i.d.\ evaluation rows.

\section{Counterexamples and invalid stronger claims}

\begin{longtable}{L{0.20\textwidth}L{0.34\textwidth}L{0.36\textwidth}}
\toprule
\PidTableHeaderRow
Lens & Construction & Consequence \\
\midrule
\endfirsthead
\toprule
\PidTableHeaderRow
Lens & Construction & Consequence \\
\midrule
\endhead
Support face & Add a new cell of mass $\varepsilon\downarrow0$. & A local informative term can grow
as $-\log\varepsilon$; uniform linear averaged bounds fail without common support. \\
Probability & Draw $C\sim\operatorname{Bernoulli}(1/2)$ and set $Z_j=C$ for every $j$. & The process
is stationary but not ergodic; a path sees only one state. \\
Rare state & Give one state mass $\varepsilon$. & It is absent after $n$ rows with probability
$(1-\varepsilon)^n$; no uniform deterministic discovery time exists. \\
$\imin$ tie & Cross a law where two specific-information values exchange order. & The minimum is
continuous but need not be differentiable; a generic delta-method normal limit does not follow. \\
Transform & Fit finite min/max ranges, then evaluate an independent standard-normal held-out law
with an error policy. & Every finite range leaves positive out-of-range mass; infinite-prefix
success has probability zero for each positive-failure training outcome. \\
Mixture & Randomly choose a correlated or anticorrelated binary push-forward law. & Each conditional
law has mutual information $\log2$, but their equal mixture is independent and has mutual
information zero. \\
Indexing & Append a rare realization that is lexicographically first. & Returned vector positions
shift; pointwise convergence requires realization keys. \\
Numerics & Exceed exact binary64 integer counts or change the platform logarithm. & The exact-real
theorem does not become a portable executable theorem. \\
Downstream & Reuse overlapping windows, drift the law, or select complete samples after inspecting
them. & The i.i.d.\ proof and its envelope do not apply. \\
\bottomrule
\end{longtable}

	The following stronger claims are invalid and are retained for audit:
	\begin{enumerate}
	  \item Strict stationarity is sufficient without ergodicity.  The latent-constant construction
	    in the table disproves it; ergodicity or another pathwise frequency premise is required.
	  \item All-unique samples identify population support.  A finite sample cannot prove that no
	    unobserved population cell has positive mass.
	  \item Pointwise atoms converge by returned-vector index.  A lexicographically earlier state
	    shifts later positions, so convergence is keyed by the realization.
	  \item One deterministic support-discovery time works for every finite law.  A cell of
	    arbitrarily small positive mass rules out a uniform time.
	  \item Complete-success conditioning preserves the original error-policy estimand.  It does
	    not; explicit rejection sampling can define a different conditional-law estimand only under
	    its own declared sampling contract.
	  \item The conditional fitted-transform limit equals the PID of the unconditional mixture.
	    Nonlinearity of PID in the law rules this out in general.
	  \item Foldwise PMFs or atoms can be pooled when each fold uses a different map.  No
	    triangular-array or cross-fit theorem establishing that operation is proved here.
	  \item The Lean artifact proves the stochastic theorem or the Rust implementation.  It proves
	    only the deterministic exact-real lemmas listed in Section~\ref{sec:evidence-boundary}.
	  \item A fixture tolerance is a portable binary64 error theorem.  The tolerance applies only to
	    the committed cases and arithmetic path.
	  \item Unrestricted exact sign and tie evaluation is always practical after denominator
	    clearing.  A rational linear combination of logarithms can reduce in principle to an
	    integer-product comparison, but the required powers can be too large for an unrestricted
	    practical implementation; no such API exists here.
	  \item The all-prefix i.i.d.\ bound covers sliding or dependent windows.  It does not; a valid
	    dependence theorem and its premises are required.
	  \item Continuity at an $\imin$ tie implies a normal limit.  Continuity alone supplies neither
	    differentiability nor a central limit theorem.
	\end{enumerate}

\section{Formal and executable evidence boundary}\label{sec:evidence-boundary}

The Lean 4.32 root project's deterministic fixed-support core proves these exact-real lemmas:
\begin{itemize}
  \item eventual positivity at a positive limit;
  \item convergence of a fixed finite event mass;
  \item positivity of a finite sum of nonnegative cell masses when the event contains a positive
    cell;
  \item sequential-limit preservation by $\log$ and $-\log$ at positive limits, and by log ratios
    at positive limits;
  \item sequential-limit preservation by $-x\log x$, including at zero;
  \item preservation of convergence by fixed finite linear and weighted sums;
  \item preservation of convergence by a finite nonempty minimum, including ties; and
  \item preservation by a finite weighted sum of finite positive-limit log-ratio combinations.
\end{itemize}
Separate modules imported by the same root project formalize the finite heterogeneous keyed Sx
event layer: source, target, and target-restricted events; their equivalence-class covers; anchor
and intersection properties; positive mass under positive anchor mass; and the law-independent
event map.  A further imported module proves a generic finite equivalence-union fractional-cover
bound and its source, target-restricted, and target-event corollaries.  The
\texttt{TwoSourceCountEventBridge.lean} module additionally defines an exact empirical law from a
supplied natural count function with positive total and proves the event-count, positivity, local
logarithm, and positive-support average identities for the source-one, source-two, joint-source,
and redundancy signed-net cumulative nodes.

The \texttt{TwoSourceMobiusAtomBridge.lean} module then fixes the categorical two-source
sub-scope's four cumulative nodes, four concrete M{\"o}bius atoms, three components, and
24-coordinate order.  It proves inverse M{\"o}bius/zeta algebra, commutation with empirical
averaging, all-coordinate supplied-count equalities, exact rational and real products, and
scaled-log sign and zero equivalences.  The paper-facing formulas are a reviewed repository
transcription, not a formal publication-to-Lean correspondence theorem.

Lean does not connect bytes or rows to that count function, encode stochastic limit theorems, the
published component-atom nonnegativity theorem, $\imin$, Rust or floating-point refinement,
certifier/parser execution, the support-change transfer between laws, a higher-source lattice, or
statistical, population, calibration, or consumer validity.  The checker
binds the complete Lake manifest and
all nine package revisions.  It verifies each dependency checkout's root, revision, origin, and
clean status with global and system Git configuration and Git environment routing disabled.  It
retains the checkout's local configuration to verify the recorded origin.  The checked project
sources cannot contain
the tokens \texttt{admit}, \texttt{axiom}, \texttt{constant}, \texttt{native\_decide},
\texttt{sorry}, or \texttt{sorryAx}.  The checker builds the project and replays its declarations
with Lean's bundled kernel checker.  It enforces an exact ordered inventory of all 339 source-level
declarations across the eight imported modules and audits the permitted axiom basis of all 246
named source theorems.  The complete two-source count/event and count-to-atom bridges are
separately SHA-256 bound.  Two separately digest-pinned semantic contracts fix paper-facing
event/count facts and an asymmetric exact 24-coordinate witness; compiled examples are not counted
as named-theorem axiom audits.  The self-test uses baseline-first static and isolated Lean changes
under normal and optimized Python to exercise imports, theorem strength, order, components,
counts, weights, products, signs, scope, and the semantic contracts.  CI runs the same formal
checker.  This does not enlarge the theorem boundary beyond supplied exact counts and the fixed
two-source categorical surface.

	The independent generator
	\nolinkurl{scripts/generate-finite-alphabet-plugin-oracle.py} uses 100-digit standard-library
	Decimal arithmetic and direct definitions; it imports no \texttt{pid-rs} code.  Its digest-bound
	fixture covers all $4$, $18$, and $166$ averaged SxPID coordinates in the listed two-, three-,
	and four-source tables, informative, misinformative, and signed-net values, realization-key
	changes after a late rare state, listed two- and three-source $\imin$ tables, and left/tie/right
	minimizer crossings.  The companion Rust test, rather than the JSON fixture, separately checks
	bit-identical outputs from fitted-quantizer wrappers and direct calls on the same categorical
	matrices, and it checks pointwise omission of an absent realization on the listed support face.
	The committed Decimal comparisons use the
	absolute envelope
\begin{equation}
  64\,\texttt{f64::EPSILON}\ \text{nats}.
\end{equation}
This is bounded implementation evidence.  It is not the convergence proof, external review,
population validation, or a global floating-point bound.

\section{Downstream contract boundary}

\begin{longtable}{L{0.15\textwidth}L{0.38\textwidth}L{0.37\textwidth}}
\toprule
\PidTableHeaderRow
Consumer & Covered path & Uncovered path \\
\midrule
\endfirsthead
\toprule
\PidTableHeaderRow
Consumer & Covered path & Uncovered path \\
\midrule
\endhead
Prisoma & Use conditionally i.i.d.\ cases or episodes from one common conditional law.  Use one
evaluation row per unit, or separately establish the row law.  Fit and freeze transforms on an
independent partition. & Independent non-identical units, dependent rows within an episode,
same-row fitting, pooling outputs from different maps, or other dependence without a theorem. \\
Galadriel & Use direct categorical SxPID on i.i.d.\ rows with a fixed alphabet. & Continuous KSG on
dependent temporal windows, same-window scaling, added-noise studies, circular delete-block work,
or complete-sample rejection after a degeneracy check. \\
Haldir & Use i.i.d.\ or strictly stationary and ergodic fixed-law rows for finite-alphabet entropy,
co-information, O-information, SxPID, and positive-denominator ratio limits. & A fixed one-time
marginal without ergodicity, drift, CUSUM, alarms, overlapping windows, post-alarm estimation, or
calibration. \\
Crebain & Use stable categorical and fitted-transform contracts after an explicit dependency and
input-contract review. & No compatibility or adapter claim follows from this theorem. \\
\bottomrule
\end{longtable}

No downstream repository is a runtime dependency of \texttt{pid-rs}.  This audit states
mathematical preconditions.  It does not claim downstream qualification.

\section{Reproduction}

Run these commands from the repository root:
\begin{verbatim}
python3 scripts/generate-finite-alphabet-plugin-oracle.py
cargo test --locked -p pid-core --test finite_alphabet_plugin_oracle
python3 scripts/check-lean-finite-convergence.py
scripts/check-finite-alphabet-convergence-pdf.sh
python3 scripts/check-method-catalog.py
\end{verbatim}

\section{Correction ledger}

The review process corrected the following defects in the current reviewed draft:
\begin{enumerate}
  \item A local-bound draft reused $p_{\min}$ from a law different from the local reference law.
    The final statement sets $p=P$.
  \item A Shannon draft substituted joint variation into a nonmonotone marginal
    Fannes--Audenaert envelope.  Equation~\eqref{eq:fannes-marginal} now clips at $1-1/K_A$.
  \item A transform draft did not require the random map to be training-measurable.  A map such as
    $Q(\omega,x)=x\mathbin{\mathrm{xor}}W_1(\omega)$ can otherwise destroy conditional
    independence.
  \item A theorem draft omitted the support-stabilized tail and allowed an empty antichain.
  \item A stopping statement omitted the condition that support conclusions need an index at
    least $N_0$.
  \item A failure statement treated a random $\eta$ as positive almost surely.  The correct
    unconditional probability is $\Prb(\eta>0)$.
  \item A rejection statement did not distinguish complete-prefix conditioning from a newly
    declared rejection-sampling estimand.
  \item A status statement incorrectly made strict threshold separation necessary for every
    eventual \texttt{Defined} sequence.  At threshold zero, the nested count path
    $(n_{00},n_{01},n_{10},n_{11})=(2,0,0,2)$ followed by repeated additions
    $01,10,00,11$ converges to independence while empirical mutual information stays positive.
  \item A downstream draft omitted common-law and ergodicity conditions.
  \item An evidence draft attributed quantizer inputs to the Decimal JSON\@.  They occur only in the
    companion Rust test.
  \item A valid first averaged-net bound used $2L$.  Applying the net cumulative range directly
    gives the tighter $L$ term in Theorem~\ref{thm:sx-local}.
  \item An intermediate entropy draft asserted $p_{\min}\le 1/K$.  This is false when $K$ is the
    declared alphabet size and $P$ has sparse support; one-cell support has $p_{\min}=1$.  No
    accepted bound uses this claim.  Alphabet size and supported-cell mass remain separate.
\end{enumerate}

\section{Open Wibral-lineage work}

The next research priority is shared-exclusions PID, not an extension of $\imin$.  Open work is:
\begin{enumerate}
  \item a dependence-aware SxPID result for generic mixing or temporal processes that do not admit
    the explicit independence coloring in the separate result under a dependency coloring;
  \item a sequential contract that separates fixed-law monitoring from drift and post-alarm
    estimation;
  \item a triangular-array or cross-fit theorem for changing fitted maps;
  \item practical certified sign and tie evaluation for empirical Sx expressions;
  \item a deductive refinement to binary64 or interval arithmetic;
  \item a Lean encoding of the empirical-law stochastic step and complete three-and-higher-source
    SxPID lattice beyond the fixed two-source supplied-count bridge, including its M{\"o}bius,
    component, cumulative, atom, and averaging composition; and
  \item mathematically valid multivariate and mixed-dimensional extensions needed by ecosystem
    consumers.
\end{enumerate}

\phantomsection
\addcontentsline{toc}{section}{References}
\begingroup
\small
\begin{thebibliography}{99}

\bibitem{audenaert2007}
K. M. R. Audenaert.
\newblock A sharp continuity estimate for the von Neumann entropy.
\newblock \emph{Journal of Physics A}, 40(28), 2007.
\newblock \url{https://doi.org/10.1088/1751-8113/40/28/S18}.

	\bibitem{bell2003}
	A. J. Bell.
	\newblock The Co-Information Lattice.
	\newblock In \emph{Proceedings of ICA 2003}, 2003.
	\newblock \url{https://www.rd.ntt/cs/team_project/icl/signal/ica2003/cdrom/data/0187.pdf}.

\bibitem{grayneuhoff1998}
R. M. Gray and D. L. Neuhoff.
\newblock Quantization.
\newblock \emph{IEEE Transactions on Information Theory}, 44(6), 1998.
\newblock \url{https://doi.org/10.1109/18.720541}.

\bibitem{gutknecht2021}
A. J. Gutknecht, M. Wibral, and A. Makkeh.
\newblock Bits and Pieces: Understanding Information Decomposition from Part-Whole Relationships
and Formal Logic.
\newblock \emph{Proceedings of the Royal Society A}, 477, 2021.
\newblock \url{https://doi.org/10.1098/rspa.2021.0110}.

\bibitem{gutknecht2025}
A. J. Gutknecht, F. E. Rosas, D. A. Ehrlich, A. Makkeh, P. A. M. Mediano, and M. Wibral.
\newblock Shannon Invariants: A Scalable Approach to Information Decomposition.
\newblock 2025.
\newblock \url{https://doi.org/10.48550/arXiv.2504.15779}.

\bibitem{hoeffding1963}
W. Hoeffding.
\newblock Probability Inequalities for Sums of Bounded Random Variables.
\newblock \emph{Journal of the American Statistical Association}, 58(301), 1963.
\newblock \url{https://doi.org/10.1080/01621459.1963.10500830}.

\bibitem{makkeh2021}
A. Makkeh, A. J. Gutknecht, and M. Wibral.
\newblock Introducing a Differentiable Measure of Pointwise Shared Information.
\newblock \emph{Physical Review E}, 103, 032149, 2021.
\newblock \url{https://doi.org/10.1103/PhysRevE.103.032149}.

\bibitem{mcgill1954}
W. J. McGill.
\newblock Multivariate Information Transmission.
\newblock \emph{Psychometrika}, 19, 1954.
\newblock \url{https://doi.org/10.1007/BF02289159}.

\bibitem{rosas2019}
F. E. Rosas, P. A. M. Mediano, M. Gastpar, and H. J. Jensen.
\newblock Quantifying High-Order Interdependencies via Multivariate Extensions of the Mutual
Information.
\newblock \emph{Physical Review E}, 100, 032305, 2019.
\newblock \url{https://doi.org/10.1103/PhysRevE.100.032305}.

\bibitem{shannon1948}
C. E. Shannon.
\newblock A Mathematical Theory of Communication.
\newblock \emph{Bell System Technical Journal}, 27, 1948.
\newblock \url{https://doi.org/10.1002/j.1538-7305.1948.tb01338.x}.

\bibitem{williamsbeer2010}
P. L. Williams and R. D. Beer.
\newblock Nonnegative Decomposition of Multivariate Information.
\newblock 2010.
\newblock \url{https://arxiv.org/abs/1004.2515}.

\end{thebibliography}
\endgroup

\end{document}
